PostgreSQL JSONB-модель строится как частично нормализованная реляционная
схема. Основные сущности ER-модели сохраняются в отдельных таблицах:
\texttt{area}, \texttt{zone}, \texttt{camera}, \texttt{event} и
\texttt{camera\_telemetry}. Связь многие-ко-многим между событиями и
камерами реализуется промежуточной таблицей \texttt{event\_camera}. Для
вариативных атрибутов используется тип \texttt{jsonb}, поддерживаемый
PostgreSQL \cite{postgresql_docs}.

\subsubsection{Ключи и связи}

В таблицах используются суррогатные первичные ключи, удобные для физического
хранения и ссылочной целостности. Предметные идентификаторы ER-модели
сохраняются как альтернативные ключи:

\begin{itemize}
  \item \texttt{area.area\_code} --- код территории;
  \item пара \texttt{zone.area\_id}, \texttt{zone.zone\_code} --- код зоны
        внутри территории;
  \item \texttt{camera.serial\_number} --- серийный номер камеры;
  \item пара \texttt{event.zone\_id}, \texttt{event.event\_number} --- номер
        события внутри зоны;
  \item пара \texttt{camera\_telemetry.camera\_id},
        \texttt{camera\_telemetry.recorded\_at} --- запись телеметрии камеры.
\end{itemize}

Связи один-ко-многим реализуются внешними ключами. Связь
\textbf{Camera}--\textbf{Event} реализуется таблицей
\texttt{event\_camera}, содержащей идентификатор события и идентификатор
камеры.

\subsubsection{Использование JSONB}

Особенность данной модели заключается в том, что составные и вариативные
атрибуты не раскрываются в отдельные отношения, а хранятся в полях типа
\texttt{jsonb}. В модели используются следующие JSONB-поля:

\begin{itemize}
  \item \texttt{camera.position} --- положение и ориентация камеры;
  \item \texttt{camera.settings} --- настройки видеопотока и аналитики;
  \item \texttt{event.payload} --- параметры события, зависящие от его типа;
  \item \texttt{camera\_telemetry.metrics} --- технические показатели камеры.
\end{itemize}

Обнаруженные объекты не выделяются в отдельную таблицу этой модели. Они
хранятся внутри \texttt{event.payload} как массив объектов события. 
Стабильная часть модели остается реляционной, а данные с переменной 
структурой хранятся как документы внутри PostgreSQL.

\subsubsection{Особенности модели}

PostgreSQL JSONB-модель позволяет оценить компромисс между реляционным
хранением и документным представлением вложенных данных. При записи события
не требуется раскладывать все параметры события и обнаруженные объекты по
множеству таблиц, однако запросы к внутренним полям \texttt{jsonb} требуют
специальных индексов и отличаются от обычных реляционных соединений.

\image{rel_jsonb}{Реляционная JSONB-модель}{0.8}
\FloatBarrier
